\section{هفته‌ی 7: ارزیابی مدل‌ها 2: معیارهای اطلاعاتی، بازیابی پارامتر و
بازیابی مدل}

\begin{quote}
\textbf{فصل 3: ارزیابی مدل‌های شناختی} · درس 1 ساعت · \lr{lab} 2 ساعت
\textbf{داده‌ی اصلی:} داده‌ها و مدل‌های فراموشی هفته‌های 2\lr{--}6، به‌علاوه‌ی
مطالعات شبیه‌سازی بزرگ با همان طرح آزمایش.
\end{quote}

\begin{center}\rule{0.5\linewidth}{0.5pt}\end{center}

\subsection{واژه‌های کلیدی}

واژه‌های مهم این درس با معادل فارسی‌شان (حتی اگر در درس‌های قبل آمده
باشند).

{\def\LTcaptype{none} % do not increment counter
\begin{longtable}[]{@{}
  >{\raggedright\arraybackslash}p{(\linewidth - 4\tabcolsep) * \real{0.3333}}
  >{\raggedright\arraybackslash}p{(\linewidth - 4\tabcolsep) * \real{0.3333}}
  >{\raggedright\arraybackslash}p{(\linewidth - 4\tabcolsep) * \real{0.3333}}@{}}
\toprule\noalign{}
\begin{minipage}[b]{\linewidth}\raggedright
\lr{English}
\end{minipage} & \begin{minipage}[b]{\linewidth}\raggedright
فارسی
\end{minipage} & \begin{minipage}[b]{\linewidth}\raggedright
توضیح کوتاه
\end{minipage} \\
\midrule\noalign{}
\endhead
\bottomrule\noalign{}
\endlastfoot
\lr{criterion (c)} & ملاک & آستانه‌ی تصمیم؛ نشان‌دهنده‌ی سوگیری پاسخ \\
\lr{parameter recovery} & بازیابی پارامتر & آیا برازش، پارامترهای معلوم
شبیه‌سازی را برمی‌گرداند \\
\lr{parameter} & پارامتر & کمیت قابل تنظیم مدل که بین افراد یا شرایط فرق
می‌کند \\
\lr{model recovery} & بازیابی مدل & آیا مقایسه‌ی مدل، مدل مولد را درست
انتخاب می‌کند \\
\lr{forgetting curve} & منحنی فراموشی & کاهش یادآوری با گذشت زمان \\
\lr{simulation} & شبیه‌سازی & تولید داده‌ی مصنوعی با اجرای مدل \\
\lr{research design} & طرح آزمایش & شرایط، تعداد کوشش‌ها و نمونه‌ها \\
\lr{likelihood-ratio test (LRT)} & آزمون نسبت درست‌نمایی & مقایسه‌ی دو مدل
تودرتو \\
\lr{likelihood} & درست‌نمایی & احتمال داده به‌عنوان تابعی از پارامترها \\
\lr{nested models} & مدل‌های تودرتو \lr{(nested)} & یک مدل حالت خاص مدل
دیگر است \\
\lr{AIC} / \lr{BIC} & معیار اطلاعاتی آکائیکه / بیزی & برازش به‌علاوه‌ی
جریمه برای تعداد پارامتر \\
\lr{response bias} & سوگیری پاسخ & تمایل به یک پاسخ، مستقل از محرک \\
\lr{confusion matrix} & ماتریس درهم‌ریختگی \lr{(confusion matrix)} &
\lr{P(}مدل انتخاب‌شده \\
\lr{inversion matrix} & ماتریس معکوس \lr{(inversion matrix)} &
\lr{P(}مدل مولد \\
\lr{fit} / \lr{fitting} & برازش & تنظیم پارامترها تا مدل با داده جور
شود \\
\lr{exponential} / \lr{power function} & تابع نمایی / توانی & دو شکل
رایج برای منحنی فراموشی و یادگیری \\
\lr{identifiability} & شناسایی‌پذیری & آیا داده پارامترها را به‌طور یکتا
تعیین می‌کند \\
\lr{parameter trade-off} & بده‌بستان پارامترها & تغییر یک پارامتر که با
پارامتر دیگر جبران می‌شود \\
\lr{log-likelihood} / \lr{NLL} & لگاریتم درست‌نمایی / منفی آن & مقداری که
بهینه‌سازها کمینه یا بیشینه می‌کنند \\
\lr{Akaike weights} & وزن‌های آکائیکه & پشتیبانی نسبی از هر مدل \\
\lr{cross-validation (CV)} & اعتبارسنجی متقابل \lr{(cross-validation)} &
برازش روی بخشی از داده و آزمون روی بخش کنارگذاشته \\
\lr{trial} & کوشش \lr{(trial)} & یک بار انجام تکلیف توسط شرکت‌کننده \\
\lr{model mimicry} & تقلید مدل & تولید داده‌ای که شبیه داده‌ی مدل دیگر
است \\
\lr{in-context learning} & یادگیری درون‌متنی & یادگیری از مثال‌های داخل
\lr{prompt} \\
\end{longtable}
}

\begin{center}\rule{0.5\linewidth}{0.5pt}\end{center}

\subsection{اهداف یادگیری}

در پایان این هفته، دانشجویان می‌توانند:

\begin{enumerate}
\def\labelenumi{\arabic{enumi}.}
\tightlist
\item
  آزمون نسبت درست‌نمایی را برای مدل‌های \lr{nested} \textbf{انجام دهند} و
  \textbf{توضیح دهند} چرا نمی‌تواند مدل‌های \lr{non-nested} را مقایسه کند.
\item
  \lr{AIC}، \lr{AICc}، \lr{BIC} و وزن‌های آکائیکه/\lr{BIC} را
  \textbf{حساب و تفسیر کنند} و اهداف متفاوت \lr{AIC} و \lr{BIC} را
  \textbf{توضیح دهند}.
\item
  یک مطالعه‌ی بازیابی پارامتر را \textbf{طراحی و اجرا کنند} و پارامترهای
  سوگیر، پرنویز یا دارای \lr{trade-off} را \textbf{تشخیص دهند}.
\item
  یک مطالعه‌ی بازیابی مدل را \textbf{طراحی و اجرا کنند} و ماتریس‌های
  درهم‌ریختگی و معکوس را \textbf{تفسیر کنند}.
\item
  \textbf{ارزیابی کنند} که آیا یک طرح آزمایش و روش مقایسه‌ی مشخص می‌تواند
  از نتیجه‌گیری‌هایی که پژوهشگر می‌خواهد بگیرد پشتیبانی کند.
\end{enumerate}

\begin{center}\rule{0.5\linewidth}{0.5pt}\end{center}

\subsection{مقدمه}

\lr{cross-validation} (هفته‌ی 6) کلی ولی پرهزینه است: هر بخش یعنی یک
برازش دوباره، و برازش دوباره‌ی مدل‌های پیچیده با \lr{multi-start} ممکن است
ساعت‌ها طول بکشد. \textbf{معیارهای اطلاعاتی} عملکرد خارج از نمونه را از
یک برازش \emph{واحد} تخمین می‌زنند؛ با افزودن یک جریمه‌ی پیچیدگی به برازش
آموزش.

ولی سؤال بنیادی‌تری هست که هیچ معیاری به‌تنهایی به آن جواب نمی‌دهد:
\textbf{آیا روش ما اصلاً با این طرح آزمایش می‌تواند کار کند؟} فرض کنید
نتیجه بگیریم «شرکت‌کننده‌ها با یک قانون توانی با نمای زوال 0.42 فراموش
می‌کنند». برای اینکه این حرف معنایی داشته باشد دو چیز باید درست باشد:

\begin{enumerate}
\def\labelenumi{\arabic{enumi}.}
\tightlist
\item
  اگر داده‌ها از یک مدل توانی با \(d = 0.42\) آمده بودند، رویه‌ی برازش ما
  چیزی نزدیک به 0.42 برمی‌گرداند (\textbf{بازیابی پارامتر}).
\item
  اگر داده‌ها از یک مدل نمایی آمده بودند، روش مقایسه‌ی ما می‌گفت «نمایی»،
  نه «توانی» (\textbf{بازیابی مدل}).
\end{enumerate}

هر دو را می‌شود با شبیه‌سازی بررسی کرد، چون برای داده‌ی شبیه‌سازی‌شده حقیقت
را می‌دانیم. راهنماهای امروزی مدل‌سازی محاسباتی این بررسی‌ها را ضروری
می‌دانند؛ نادیده گرفتنشان مثل اجرای آزمایش بدون وارسی سالم بودن تجهیزات
است. این هفته بخش «شبیه‌سازی \ensuremath{\leftarrow} برازش \ensuremath{\leftarrow} ارزیابی \ensuremath{\leftarrow} \textbf{بازیابی}» مسیر
مدل‌سازی ما را کامل می‌کند.

\paragraph{این موضوع کجا به کار می‌رود؟}

\begin{itemize}
\tightlist
\item
  \textbf{جدول‌های مقایسه‌ی مدل با \lr{AIC/BIC}} تقریباً در هر مقاله‌ی
  مدل‌سازی در روان‌شناسی، علوم اعصاب، بوم‌شناسی و اقتصاد دیده می‌شوند.
\item
  \textbf{روان‌پزشکی محاسباتی:} پیش از ادعای اینکه بیماران نرخ یادگیری
  پایین‌تری دارند، پژوهشگرها باید با بازیابی پارامتر نشان دهند که نرخ
  یادگیری در طرح آزمایش آن‌ها قابل برآورد است. داوران مقاله‌ها روزبه‌روز
  بیشتر این را می‌خواهند.
\item
  \textbf{طراحی آزمایش:} شبیه‌سازی‌های بازیابی \emph{پیش از} جمع‌آوری داده
  اجرا می‌شوند تا تعداد \lr{trial}ها و شرایط انتخاب شود (نسخه‌ی مدل‌سازی
  تحلیل توان آماری یا \lr{power analysis).}
\item
  \textbf{هر مطالعه‌ای که می‌گوید «مدل \lr{X} داده را بهتر از همه توضیح
  می‌دهد»} باید بازیابی مدل را گزارش کند؛ وگرنه ممکن است برنده فقط سوگیری
  معیار مقایسه را نشان دهد.
\end{itemize}

\begin{center}\rule{0.5\linewidth}{0.5pt}\end{center}

\subsection{جزوه‌ی درس}

\paragraph{به زبان ساده}

\lr{cross-validation} صادق ولی کند است. \textbf{معیارهای اطلاعاتی}
\lr{(AIC} و \lr{BIC)} فرمول‌های سریعی‌اند که برازش هر مدل را می‌گیرند و
برای هر پارامتر یک جریمه کم می‌کنند؛ پس مدلی برنده می‌شود که پارامترهای
اضافه‌اش برازش را بیش از «هزینه‌شان» بهتر کنند. هدف \lr{AIC} بهترین
پیش‌بینی‌کننده و هدف \lr{BIC} مدل درست است، و جریمه‌ی \lr{BIC} بزرگ‌تر است.
نیمه‌ی دوم هفته که مهم‌تر است می‌پرسد آیا روش‌های ما اصلاً کار می‌کنند. آن‌ها
را روی داده‌ی ساختگی امتحان می‌کنیم که در آن حقیقت را \emph{می‌دانیم}.
\textbf{بازیابی پارامتر:} داده را با پارامترهای معلوم شبیه‌سازی کن، برازش
کن، و بررسی کن که مقدارهای برازش‌شده نزدیک مقدارهای واقعی برگردند.
\textbf{بازیابی مدل:} داده را از هر مدل نامزد شبیه‌سازی کن و بررسی کن که
روش مقایسه همان مدلی را انتخاب کند که داده را تولید کرده. اگر روشی
نتواند حقیقت را در داده‌ی شبیه‌سازی‌شده پیدا کند، جوابش روی داده‌ی واقعی
بی‌معناست.

\textbf{زمان‌بندی پیشنهادی (60 دقیقه):} آزمون نسبت درست‌نمایی 8 · \lr{AIC}
و \lr{BIC} 20 · بازیابی پارامتر 14 · بازیابی مدل 14 ·
\lr{identifiability} 4.

\subsubsection{\texorpdfstring{1. آزمون نسبت درست‌نمایی (مدل‌های
\lr{nested)}}{1. آزمون نسبت درست‌نمایی (مدل‌های nested)}}

مدل \(\mathcal{M}_0\) در \(\mathcal{M}_1\) \textbf{\lr{nested}} است اگر
\(\mathcal{M}_0\) حالت خاصی از \(\mathcal{M}_1\) باشد که با ثابت کردن
بعضی پارامترها به دست می‌آید. مثال: مدل نمایی همان مدل نمایی با مجانب با
\(s = 0\) (مجانب در 0) است. تحت فرض صفر که \(\mathcal{M}_0\) درست است،

\[
\Lambda = 2\big[\ell_1(\hat{\boldsymbol\theta}_1) - \ell_0(\hat{\boldsymbol\theta}_0)\big] \;\dot\sim\; \chi^2_{\Delta k},
\]

که در آن \(\Delta k\) تفاوت تعداد پارامترهای آزاد است. \(\Lambda\) بزرگ
یعنی پارامترهای اضافه برازش را بیشتر از آنچه شانس توضیح می‌دهد بهتر
کرده‌اند.

\textbf{محدودیت‌ها:} فقط برای مدل‌های \textbf{\lr{nested}} کار می‌کند
(توانی در برابر نمایی را نمی‌شود این‌طور آزمون کرد)؛ آزمون فرض صفر است
(می‌تواند مدل ساده‌تر را رد کند ولی هرگز از آن حمایت نمی‌کند)؛ و وقتی مقدار
ثابت‌شده‌ی پارامتر روی مرز فضای پارامتر است (مثل \(s = 0\)) تقریبی است.

\begin{quote}
\textbf{پیشرفته: فرض صفر روی مرز.} وقتی مقدار صفر روی مرز است (واریانس
0، مجانب 0، وزن آمیخته‌ی 0)، توزیع \(\Lambda\) تحت فرض صفر آمیزه‌ای 50:50
از \(\chi^2_0\) (یک جرم نقطه‌ای در 0) و \(\chi^2_1\) است. در این حالت
استفاده از \(\chi^2_1\) محافظه‌کارانه است: مقدار \(p\) درست نصف مقدار
ساده است.
\end{quote}

\subsubsection{2. معیارهای اطلاعاتی}

معیارهای اطلاعاتی بین برازش و پیچیدگی \lr{trade-off} برقرار می‌کنند. با
لگاریتم درست‌نمایی بیشینه \(\hat\ell\)، \(k\) پارامتر آزاد و \(N\)
مشاهده:

\[
\mathrm{AIC} = -2\hat\ell + 2k, \qquad
\mathrm{AICc} = \mathrm{AIC} + \frac{2k(k+1)}{N - k - 1}, \qquad
\mathrm{BIC} = -2\hat\ell + k \ln N.
\]

\textbf{کمتر بهتر است.} فقط \textbf{تفاوت‌ها} بین مدل‌هایی که روی داده‌ی
یکسان برازش شده‌اند معنادارند.

{\def\LTcaptype{none} % do not increment counter
\begin{longtable}[]{@{}
  >{\raggedright\arraybackslash}p{(\linewidth - 4\tabcolsep) * \real{0.3333}}
  >{\raggedright\arraybackslash}p{(\linewidth - 4\tabcolsep) * \real{0.3333}}
  >{\raggedright\arraybackslash}p{(\linewidth - 4\tabcolsep) * \real{0.3333}}@{}}
\toprule\noalign{}
\begin{minipage}[b]{\linewidth}\raggedright
\end{minipage} & \begin{minipage}[b]{\linewidth}\raggedright
\lr{AIC} (و \lr{AICc)}
\end{minipage} & \begin{minipage}[b]{\linewidth}\raggedright
\lr{BIC}
\end{minipage} \\
\midrule\noalign{}
\endhead
\bottomrule\noalign{}
\endlastfoot
هدف & بهترین \textbf{پیش‌بینی} داده‌ی جدید (کمینه کردن واگرایی
\lr{Kullback\lr{--}Leibler} مورد انتظار نسبت به حقیقت) & انتخاب مدل
\textbf{واقعی} اگر میان نامزدها باشد (سازگاری)؛ تقریبی از \lr{Bayes
factor} \\
جریمه‌ی هر پارامتر & 2 & \(\ln N\) (از \(N \ge 8\) بیشتر از 2) \\
وقتی \(N \to \infty\) & ممکن است همچنان مدل‌های بیش از حد پیچیده را
انتخاب کند & مدل واقعی را با احتمال \ensuremath{\leftarrow} 1 انتخاب می‌کند \\
تصحیح نمونه‌ی کوچک & \lr{AICc} (وقتی \(N/k < 40\)) & \lr{---} \\
\end{longtable}
}

\textbf{\(N\) چیست؟} تعداد \emph{مشاهدات مستقلی} که وارد درست‌نمایی
می‌شوند. برای داده‌ی یادآوری دوجمله‌ای این تعداد پاسخ‌های آیتم است (اینجا
\(8 \times 50 = 400\))، نه تعداد خانه‌ها (8). این انتخاب روی \lr{BIC} (از
راه \(\ln N\)) و \lr{AICc} اثر دارد. آن را گزارش کنید.

\textbf{وزن‌ها.} با \(\Delta_i = \mathrm{IC}_i - \min_j \mathrm{IC}_j\)،

\[
w_i = \frac{\exp(-\Delta_i/2)}{\sum_j \exp(-\Delta_j/2)}.
\]

\textbf{وزن‌های آکائیکه} احتمال اینکه مدل \(i\) بهترین پیش‌بینی‌کننده در
میان نامزدها باشد را تخمین می‌زنند. \textbf{وزن‌های \lr{BIC}} احتمال‌های
پسین مدل‌ها را تحت بخت‌های پیشین برابر تقریب می‌زنند (هفته‌ی 10).

\textbf{قواعد سرانگشتی برای \(\Delta\):} 0 تا 2، پشتیبانی عملاً برابر؛ 4
تا 7، پشتیبانی به‌طور قابل توجهی کمتر؛ بیشتر از 10، عملاً بدون پشتیبانی.
این‌ها قرارداد هستند، نه قانون.

\textbf{محدودیت مشترک \lr{AIC} و \lr{BIC:}} هر دو پارامترها (\(k\)) را
می‌شمارند، پس به تفاوت \textbf{انعطاف شکل تابعی} بین مدل‌هایی با \(k\)
یکسان کورند (هفته‌ی 6). \lr{cross-validation} و \lr{Bayes factor} (هفته‌ی
10) این‌طور نیستند.

\begin{quote}
\textbf{پیشرفته: جریمه‌ی \lr{AIC} از کجا می‌آید.} لگاریتم درست‌نمایی بیشینه
برآوردی خوش‌بینانه و سوگیر از لگاریتم درست‌نمایی مورد انتظار داده‌ی جدید
است، چون همان داده برای انتخاب \(\hat{\boldsymbol\theta}\) به کار رفته
است. آکائیکه نشان داد که برای مدل‌های منظم و \(N\) بزرگ، این سوگیری
تقریباً \(k\) است. تصحیحش \(-2(\hat\ell - k) = \mathrm{AIC}\) را می‌دهد.
به همین دلیل است که \lr{AIC} و \lr{LOO-CV} به‌طور مجانبی با هم توافق
دارند.
\end{quote}

\paragraph{خواندن ریاضی}

\begin{itemize}
\tightlist
\item
  \(-2\hat\ell\) بخش \textbf{عدم برازش} است (دو برابر \lr{NLL}
  کمینه‌شده)؛ \(2k\) یا \(k\ln N\) بخش \textbf{جریمه} است. \lr{AIC/BIC}
  کمتر بهتر است.
\item
  هر پارامتر اضافه باید \(-2\hat\ell\) را بیش از 2 (برای \lr{AIC)} یا
  بیش از \(\ln N\) (برای \lr{BIC)} کم کند تا ارزشش را داشته باشد. با
  \(N = 400\)، \(\ln N \approx 6\): یک پارامتر باید \lr{NLL} را حدود 3
  واحد بهتر کند تا از \lr{BIC} بگذرد، ولی برای \lr{AIC} فقط 1 واحد.
\item
  \textbf{وزن‌ها:} \(\exp(-\Delta_i/2)\) تفاوت \lr{AIC} را به درست‌نمایی
  نسبی تبدیل می‌کند. \(\Delta = 2\) می‌دهد \(e^{-1} = 0.37\)، پس مدلی که 2
  واحد \lr{AIC} بدتر است حدود یک‌سوم پشتیبانی دارد. تقسیم بر جمع، وزن‌ها
  را طوری می‌کند که جمعشان 1 شود.
\item
  \textbf{عددهای مثال حل‌شده:} \lr{AIC} مدل توانی
  \(= 2(19.830) + 2(2) = 43.66\)؛ مدل نمایی با مجانب
  \(= 2(18.919) + 2(3) = 43.84\). مدل مجانب بهتر برازش می‌شود \lr{(NLL}
  کمتر) ولی هزینه‌ی یک پارامتر بیشتر را می‌دهد، پس دو مدل تقریباً مساوی
  می‌شوند.
\end{itemize}

\subsubsection{3. بازیابی پارامتر}

\textbf{سؤال:} با این طرح آزمایش، آیا رویه‌ی برازش پارامترهایی را که داده
را تولید کرده‌اند برمی‌گرداند؟

\textbf{رویه:} 1. مجموعه‌ای از مقادیر واقعی پارامتر
\(\boldsymbol\theta^\ast\) انتخاب کنید که دامنه‌ی معقول را پوشش دهد (یک
شبکه، یا نمونه‌های تصادفی از دامنه‌های معقول، مثلاً حول برآوردهای داده‌ی
واقعی). 2. برای هر \(\boldsymbol\theta^\ast\) یک مجموعه‌داده با
\textbf{دقیقاً همان طرح واقعی} (همان شرایط، همان تعداد \lr{trial)}
شبیه‌سازی کنید. 3. مدل را با \textbf{دقیقاً همان رویه‌ی برازش واقعی} (همان
بهینه‌ساز، شروع‌ها، کران‌ها) برازش کنید. 4. \(\hat{\boldsymbol\theta}\) را
با \(\boldsymbol\theta^\ast\) مقایسه کنید: نمودار پراکنش واقعی در برابر
بازیابی‌شده؛ همبستگی؛ سوگیری (میانگین \(\hat\theta - \theta^\ast\))؛
\lr{RMSE.} 5. پارامترهای بازیابی‌شده را \textbf{در برابر هم} هم رسم کنید.
همبستگی قوی بین مثلاً \(\hat\lambda\) و \(\hat s\) که در مقادیر واقعی
نیست، یک \textbf{\lr{trade-off}} را آشکار می‌کند.

\textbf{خواندن نتایج:} نقاط روی خط همانی یعنی بازیابی خوب؛ ابری موازی با
محور افقی یعنی پارامتر با این طرح قابل بازیابی نیست؛ منحنی‌ای که در دو
انتها کج می‌شود یعنی سوگیری نزدیک کران‌ها. اگر بازیابی شکست بخورد، آن
پارامتر را تفسیر نکنید. طرح آزمایش را عوض کنید \lr{(trial} بیشتر، شرایط
دیگر) یا مدل را ساده کنید.

\subsubsection{4. بازیابی مدل}

\textbf{سؤال:} اگر مدل \(\mathcal{M}_a\) داده را تولید کرده باشد، آیا
روش مقایسه‌ی ما \(\mathcal{M}_a\) را انتخاب می‌کند؟

\textbf{رویه:} 1. برای هر مدل نامزد \(\mathcal{M}_a\): مجموعه‌داده‌های
زیادی با طرح واقعی و با پارامترهایی از دامنه‌های معقول شبیه‌سازی کنید. 2.
\textbf{همه‌ی} مدل‌های نامزد را روی هر مجموعه‌داده برازش کنید. 3. با معیار
انتخاب‌شده \lr{(AIC}، \lr{BIC}، \lr{CV}، \ensuremath{\ldots}) یک برنده انتخاب کنید. 4.
\textbf{\lr{confusion matrix}}
\(P(\text{\lr{selected}} = b \mid \text{\lr{generated}} = a)\) را جدول کنید.

یک روش خوب ماتریسی با درایه‌های قطری بزرگ می‌دهد. جرم خارج از قطر
\textbf{تقلید مدل} \lr{(model mimicry)} را نشان می‌دهد: یک مدل می‌تواند
داده‌هایی بسازد که شبیه داده‌های مدل دیگر است.

چیزی که در تفسیر داده‌ی واقعی واقعاً لازم داریم برعکس است: با فرض اینکه
مدل \(b\) برنده شد، چقدر محتمل است که \(b\) داده را تولید کرده باشد؟ این
\textbf{\lr{inversion matrix}}
\(P(\text{\lr{generated}} = a \mid \text{\lr{selected}} = b)\) است که با نرمال
کردن هر \emph{ستون} شمارش‌های درهم‌ریختگی به دست می‌آید (با فرض فراوانی
پیشین برابر برای مدل‌های تولیدکننده).

\begin{quote}
\textbf{نکته:} بازیابی مدل \textbf{سوگیری‌های معیارها} را هم آشکار می‌کند:
\lr{BIC} در \(N\) کوچک به مدل‌های ساده‌تر و \lr{AIC} به مدل‌های پیچیده‌تر
تمایل دارد. «بهترین» معیار آن است که مدل‌های نامزد شما را \emph{با طرح
آزمایش شما} بهتر بازیابی کند، و این دقیقاً همان چیزی است که این مطالعه
اندازه می‌گیرد.
\end{quote}

\paragraph{خواندن ریاضی}

\begin{itemize}
\tightlist
\item
  \textbf{\lr{confusion matrix}}
  \(P(\text{\lr{selected}} = b \mid \text{\lr{generated}} = a)\): سطرها مدل‌های
  واقعی (مولد) و ستون‌ها مدل‌های انتخاب‌شده‌اند. جمع هر سطر 1 است. هر سطر را
  این‌طور بخوانید: «وقتی حقیقت \(a\) است، چند بار هر مدل را انتخاب
  می‌کنیم؟»
\item
  \textbf{\lr{inversion matrix}}
  \(P(\text{\lr{generated}} = a \mid \text{\lr{selected}} = b)\): همان شمارش‌ها،
  ولی هر \emph{ستون} به 1 نرمال شده. هر ستون را این‌طور بخوانید: «وقتی
  \(b\) را انتخاب کردیم، چند بار واقعاً \(b\) حقیقت بود؟» این همان سؤالی
  است که برای داده‌ی واقعی اهمیت دارد، جایی که فقط برنده را می‌بینیم.
\item
  \textbf{مثال کوچک.} از هر کدام از دو مدل 100 مجموعه داده. از \lr{A:}
  90 بار \lr{A} و 10 بار \lr{B} انتخاب می‌شود. از \lr{B:} 40 بار \lr{A} و
  60 بار \lr{B.} سطر \lr{A} در \lr{confusion matrix} برابر (0.90, 0.10)
  است. ولی از 130 مجموعه داده‌ای که \lr{A} را انتخاب کردند، فقط 90 تا از
  \lr{A} آمده بودند:
  \(P(\text{\lr{A true}} \mid \text{\lr{A selected}}) = 90/130 = 0.69\).
\end{itemize}

\subsubsection{\texorpdfstring{5.
\lr{identifiability}}{5. identifiability}}

\begin{itemize}
\tightlist
\item
  \textbf{\lr{non-identifiable} ساختاری:} مقادیر مختلف پارامتر برای هر
  داده‌ای \emph{دقیقاً} پیش‌بینی‌های یکسان می‌دهند. مثال‌ها: جابه‌جایی برچسب در
  مدل دوردی (هفته‌ی 3)، یا مکان و مقیاس در \lr{SDT} (هفته‌ی 2). راه‌حل یک
  قید یا \lr{reparameterization} است.
\item
  \textbf{\lr{non-identifiable} عملی:} پارامترها در اصل
  \lr{identifiable} هستند ولی با داده‌های موجود نه. بازیابی ضعیف است و
  برآوردها \lr{trade-off} دارند. راه‌حل طرح آزمایش بهتر یا مدل ساده‌تر
  است.
\end{itemize}

بازیابی پارامتر آزمون عملی \lr{identifiability} است.

\paragraph{برای یک مهندس}

\begin{itemize}
\tightlist
\item
  \textbf{\lr{AIC} و \lr{BIC}} همان معیارهای انتخاب مرتبه‌ی مدل‌اند که
  شاید از مدل‌سازی \lr{AR/ARMA} بشناسید؛ \textbf{\lr{BIC}} اساساً همان
  معیار \textbf{\lr{MDL}} ریسانن است. \lr{LRT} برای مدل‌های \lr{nested}
  دوباره همان \lr{GLRT} است.
\item
  \textbf{بازیابی پارامتر} یعنی آزمودن یک تخمین‌گر روی سیگنال‌های مصنوعی
  با حقیقت معلوم، با همان \lr{SNR} و طول داده‌ی واقعی، پیش از اعتماد به
  آن روی داده‌ی اندازه‌گیری‌شده. مهندس‌ها این کار را مرتب انجام می‌دهند؛ در
  روان‌شناسی اغلب حذف می‌شود، و به همین دلیل این درس روی آن اصرار دارد.
\item
  \textbf{\lr{confusion matrix}} همان ماتریس احتمال انتقال یک کانال
  گسسته است: ورودی = مدل مولد، خروجی = مدل انتخاب‌شده.
  \textbf{\lr{inversion matrix}} ماتریس احتمال‌های \emph{پسین}
  \(P(\text{\lr{input}} \mid \text{\lr{output}})\) است، یعنی چیزی که یک دیکدر
  \lr{MAP} لازم دارد، با فرض احتمال‌های پیشین برابر برای ورودی‌ها.
\item
  \textbf{\lr{non-identifiability} ساختاری} مثل یک حالت مشاهده‌ناپذیر یا
  ماتریس سیستم با رتبه‌ی ناقص است: هیچ مقداری از داده کمک نمی‌کند.
  \textbf{\lr{non-identifiability} عملی} یک مسئله‌ی بدحالت است: در اصل
  قابل شناسایی است ولی برای داده‌ی موجود نویز خیلی زیاد است.
\end{itemize}

\paragraph{برای مدرس}

\begin{itemize}
\tightlist
\item
  \textbf{بازیابی را «آزمودن ابزار» معرفی کنید.} هیچ‌کس به دماسنجی که
  هیچ‌وقت در آبی با دمای معلوم امتحان نشده اعتماد نمی‌کند. بازیابی همان
  آزمون برای یک مدل است.
\item
  \textbf{پیش از هر نظریه‌ای یک نمودار پراکندگی مقدار واقعی در برابر
  بازیابی‌شده نشان دهید} (بخش \lr{C lab):} بازیابی خوب (نقطه‌ها روی قطر)
  کنار بازیابی بد (یک ابر صاف). تصویر خودش توضیح می‌دهد.
\item
  \textbf{\lr{confusion} و \lr{inversion} را با مثال کوچک 2×2 بالا درس
  دهید}؛ دانشجوها در امتحان بیش از هر چیز دیگری در این هفته این دو را با
  هم قاطی می‌کنند. تشبیه تست پزشکی کمک می‌کند: \lr{sensitivity} یک تست با
  احتمال بیمار بودن کسی که تستش مثبت شده یکی نیست.
\item
  \textbf{سؤال‌های رایج.}

  \begin{itemize}
  \tightlist
  \item
    \emph{\lr{«AIC} یا \lr{BIC}؟»} به هدف بستگی دارد (پیش‌بینی در برابر
    پیدا کردن مدل درست) و در عمل به اینکه کدام در طرح آزمایش شما مدل‌ها
    را بهتر بازیابی می‌کند.
  \item
    \emph{«چند مجموعه داده‌ی شبیه‌سازی‌شده؟»} برای دیدن الگو 50 تا 100 برای
    هر حالت کافی است؛ برای انتشار بیشتر.
  \item
    \emph{«بازیابی شکست خورد، حالا چه؟»} آن پارامتر را تفسیر نکنید. طرح
    آزمایش را تغییر دهید یا مدل را ساده‌تر کنید.
  \end{itemize}
\end{itemize}

\begin{center}\rule{0.5\linewidth}{0.5pt}\end{center}

\subsection{مثال حل‌شده}

چهار مدل فراموشی که با \lr{MLE} روی داده‌های هفته‌ی 2 برازش شده‌اند
\lr{(NLL}های آموزش از هفته‌ی 6)، با \(N = 400\) پاسخ آیتم:

{\def\LTcaptype{none} % do not increment counter
\begin{longtable}[]{@{}
  >{\raggedright\arraybackslash}p{(\linewidth - 8\tabcolsep) * \real{0.2000}}
  >{\raggedright\arraybackslash}p{(\linewidth - 8\tabcolsep) * \real{0.2000}}
  >{\raggedright\arraybackslash}p{(\linewidth - 8\tabcolsep) * \real{0.2000}}
  >{\raggedright\arraybackslash}p{(\linewidth - 8\tabcolsep) * \real{0.2000}}
  >{\raggedright\arraybackslash}p{(\linewidth - 8\tabcolsep) * \real{0.2000}}@{}}
\toprule\noalign{}
\begin{minipage}[b]{\linewidth}\raggedright
مدل
\end{minipage} & \begin{minipage}[b]{\linewidth}\raggedright
\(k\)
\end{minipage} & \begin{minipage}[b]{\linewidth}\raggedright
\lr{NLL}
\end{minipage} & \begin{minipage}[b]{\linewidth}\raggedright
\lr{AIC} \(= 2\,\mathrm{NLL} + 2k\)
\end{minipage} & \begin{minipage}[b]{\linewidth}\raggedright
\lr{BIC} \(= 2\,\mathrm{NLL} + k\ln 400\)
\end{minipage} \\
\midrule\noalign{}
\endhead
\bottomrule\noalign{}
\endlastfoot
\lr{exp} & 2 & 29.812 & 63.62 & 71.61 \\
\lr{pow} & 2 & 19.830 & \textbf{43.66} & \textbf{51.64} \\
\lr{expasym} & 3 & 18.919 & 43.84 & 55.81 \\
\lr{two} & 4 & 18.919 & 45.84 & 61.80 \\
\end{longtable}
}

\textbf{قدم 1: \lr{AIC} برای مدل توانی.}
\(2 \times 19.830 + 2 \times 2 = 43.66\).

\textbf{قدم 2: \lr{BIC} برای مدل توانی.} \(\ln 400 = 5.991\)، پس
\(2 \times 19.830 + 2 \times 5.991 = 51.64\).

\textbf{قدم 3: وزن‌های آکائیکه.} تفاوت‌ها با بهترین \lr{AIC:}
\(\Delta = (19.96, 0, 0.18, 2.18)\). پس
\(e^{-\Delta/2} = (0.00005,\allowbreak 1,\allowbreak 0.914,\allowbreak 0.336)\) که جمعشان 2.250 است،
بنابراین

\[
w^{\text{\lr{AIC}}} = (0.000,\ 0.444,\ 0.406,\ 0.149).
\]

وزن‌های \lr{BIC} به همین روش: \((0.000,\ 0.884,\ 0.110,\ 0.005)\).

\textbf{قدم 4: آزمون نسبت درست‌نمایی}، \lr{exp nested} در \lr{expasym:}
\(\Lambda = 2(29.812 - 18.919) = 21.79\) با 1 درجه‌ی آزادی،
\(p \approx 3 \times 10^{-6}\). افزودن مجانب مدل نمایی را به‌طور معناداری
بهتر می‌کند.

\textbf{تفسیر.} هر دو معیار مدل نمایی خالص را قاطعانه رد می‌کنند.
\lr{AIC} نمی‌تواند مدل‌های توانی و مجانبی را از هم جدا کند (وزن‌های 0.44 در
برابر 0.41): از نظر \emph{پیش‌بینی} تقریباً به یک اندازه خوب‌اند، که با
نتایج \lr{CV} هفته‌ی 6 سازگار است. جریمه‌ی بزرگ‌تر \lr{BIC} مدل توانی
دوپارامتری را ترجیح می‌دهد (0.88). پارامتر اضافه‌ی مدل دوردی هیچ چیزی
نمی‌خرد: \lr{NLL} آن با \lr{expasym} یکسان است، پس جریمه را بی‌دلیل
می‌پردازد. اینکه آیا باید به این ترجیح‌ها با این طرح آزمایش \emph{اعتماد}
کنیم سؤالی برای بازیابی مدل است (بخش \lr{D lab).}

\begin{center}\rule{0.5\linewidth}{0.5pt}\end{center}

\subsection{\texorpdfstring{\lr{lab} (2 ساعت)}{lab (2 ساعت)}}

\subsubsection{اهداف}

\begin{enumerate}
\def\labelenumi{\arabic{enumi}.}
\tightlist
\item
  اجرای آزمون نسبت درست‌نمایی و محاسبه‌ی \lr{AIC}، \lr{AICc}، \lr{BIC} و
  وزن‌ها برای مدل‌های فراموشی.
\item
  اجرای یک مطالعه‌ی بازیابی پارامتر برای مدل‌های توانی و نمایی با مجانب و
  بررسی \lr{trade-off}ها.
\item
  اجرای یک مطالعه‌ی بازیابی مدل برای سه مدل با \lr{AIC} و \lr{BIC}؛ ساختن
  ماتریس‌های درهم‌ریختگی و معکوس.
\item
  \textbf{تفسیر}: با این طرح آزمایش، از داده‌های هفته‌ی 2 چه چیزی را می‌شود
  و چه چیزی را نمی‌شود نتیجه گرفت؟
\end{enumerate}

\textbf{زمان‌بندی:} بخش \lr{A} 15 دقیقه · بخش \lr{B} 15 دقیقه · بخش
\lr{C} 35 دقیقه · بخش \lr{D} 40 دقیقه · بخش \lr{E} 15 دقیقه.

\subsubsection{آماده‌سازی}

\begin{quote}
\textbf{\lr{Code} · کد:} \lr{\nolinkurl{code/week07.py}}, \lr{cell} 1 ·
سلول 1
\end{quote}

\subsubsection{\texorpdfstring{بخش \lr{A:} آزمون نسبت درست‌نمایی (15
دقیقه)}{بخش A: آزمون نسبت درست‌نمایی (15 دقیقه)}}

\begin{quote}
\textbf{\lr{Code} · کد:} \lr{\nolinkurl{code/week07.py}}, \lr{cell} 2 ·
سلول 2
\end{quote}

\textbf{خروجی مورد انتظار:}

\Needspace{7\baselineskip}
\begin{latin}
\begin{Verbatim}
NLL exp = 29.812, NLL expasym = 18.919
LR statistic = 21.79, df = 1, naive p = 3.05e-06, boundary-corrected p = 1.52e-06
\end{Verbatim}
\end{latin}

\textbf{سؤال:} چرا نمی‌توانید از همین آزمون برای مقایسه‌ی مدل‌های توانی و
نمایی استفاده کنید؟

\subsubsection{\texorpdfstring{بخش \lr{B:} معیارهای اطلاعاتی و وزن‌ها (15
دقیقه)}{بخش B: معیارهای اطلاعاتی و وزن‌ها (15 دقیقه)}}

\begin{quote}
\textbf{\lr{Code} · کد:} \lr{\nolinkurl{code/week07.py}}, \lr{cell} 3 ·
سلول 3
\end{quote}

\textbf{خروجی مورد انتظار:}

\Needspace{11\baselineskip}
\begin{latin}
\begin{Verbatim}
         k     NLL     AIC    AICc     BIC  w_AIC  w_AICc  w_BIC
model
exp      2  29.812  63.625  63.655  71.608  0.000   0.000  0.000
pow      2  19.830  43.661  43.691  51.644  0.444   0.449  0.884
expasym  3  18.919  43.837  43.898  55.812  0.406   0.405  0.110
two      4  18.919  45.837  45.939  61.803  0.150   0.146  0.006
\end{Verbatim}
\end{latin}

\textbf{بررسی حساسیت:} \lr{BIC} را با \(N = 8\) (تعداد خانه‌ها) به‌جای 400
دوباره حساب کنید. آیا رتبه‌بندی عوض می‌شود؟ کدام انتخاب \(N\) اینجا قابل
دفاع است و چرا؟

\subsubsection{\texorpdfstring{بخش \lr{C:} بازیابی پارامتر (35
دقیقه)}{بخش C: بازیابی پارامتر (35 دقیقه)}}

\begin{quote}
\textbf{\lr{Code} · کد:} \lr{\nolinkurl{code/week07.py}}, \lr{cell} 4 ·
سلول 4
\end{quote}

\textbf{خروجی مورد انتظار} (حدود 1 دقیقه):

\Needspace{17\baselineskip}
\begin{latin}
\begin{Verbatim}
pow:
   r(true, rec)   bias   RMSE
A         0.888  0.005  0.055
d         0.965  0.003  0.070

expasym:
     r(true, rec)   bias   RMSE
A           0.802  0.008  0.077
lam         0.799  0.034  0.230
s           0.918 -0.008  0.060

expasym: corr(true lam, true s) = -0.08;  corr(recovered lam, recovered s) = 0.03
\end{Verbatim}
\end{latin}

\textbf{تفسیر:} 1. با این طرح آزمایش، کدام پارامترها خوب بازیابی می‌شوند؟
تفسیر کدام را با تردید انجام می‌دهید؟ 2. نرخ زوال \(\lambda\) در
\lr{expasym} با فاصله‌ی زیاد بیشترین \lr{RMSE} را دارد (0.23 روی دامنه‌ی
0.05 تا 1). به نمودارش نگاه کنید: بازیابی برای کدام مقادیر واقعی شکست
می‌خورد و چرا؟ (راهنمایی: اگر زوال خیلی سریع باشد، بیشترِ افت پیش از دومین
فاصله‌ی آزمون اتفاق می‌افتد؛ پس داده‌ها چه شکلی پیدا می‌کنند؟) اینجا
\(\lambda\) و \(s\) بازیابی‌شده همبستگی \emph{ندارند} \lr{(r} = 0.03)، پس
این عدم دقت است نه \lr{trade-off.} این را با \lr{trade-off} آشکار
\(A\)\lr{--}\(d\) مدل توانی که در هفته‌ی 5 دیدیم مقایسه کنید. 3. چه تغییری در
طرح آزمایش بازیابی \(\lambda\) و \(s\) را بهتر می‌کند؟ یک ایده را در
تمرین 3 آزمون کنید.

\subsubsection{\texorpdfstring{بخش \lr{D:} بازیابی مدل (40
دقیقه)}{بخش D: بازیابی مدل (40 دقیقه)}}

\begin{quote}
\textbf{\lr{Code} · کد:} \lr{\nolinkurl{code/week07.py}}, \lr{cell} 5 ·
سلول 5
\end{quote}

\textbf{خروجی مورد انتظار} (حدود 1 تا 2 دقیقه):

\Needspace{26\baselineskip}
\begin{latin}
\begin{Verbatim}
AIC confusion matrix (rows = generating model, columns = selected model):
          exp   pow  expasym
exp      0.91  0.00     0.09
pow      0.03  0.82     0.15
expasym  0.06  0.25     0.69
AIC inversion matrix (rows = generating model, columns = selected model):
          exp   pow  expasym
exp      0.91  0.00     0.10
pow      0.03  0.77     0.16
expasym  0.06  0.23     0.74

BIC confusion matrix (rows = generating model, columns = selected model):
          exp   pow  expasym
exp      0.99  0.01     0.00
pow      0.05  0.94     0.01
expasym  0.09  0.58     0.33
BIC inversion matrix (rows = generating model, columns = selected model):
          exp   pow  expasym
exp      0.88  0.01     0.00
pow      0.04  0.61     0.03
expasym  0.08  0.38     0.97
\end{Verbatim}
\end{latin}

\textbf{تفسیر:} 1. کدام مدل بیشتر با کدام اشتباه گرفته می‌شود؟ (با
\lr{AIC}، 25\% مجموعه‌داده‌های \lr{expasym} برچسب توانی می‌گیرند؛ با
\lr{BIC}، 58\%.) این تقلید را بر حسب شکل منحنی‌ها توضیح دهید. 2. \lr{AIC}
و \lr{BIC} را مقایسه کنید. کدام مدل‌های ساده‌تر را بهتر بازیابی می‌کند و
کدام مدل پیچیده‌تر را؟ 3. در بخش \lr{B}، \lr{BIC} برای داده‌ی واقعی مدل
توانی را انتخاب کرد. با استفاده از \textbf{\lr{inversion matrix}}
\lr{BIC}، چقدر باید مطمئن باشیم که داده‌ها را یک مدل توانی تولید کرده
است؟ (فقط 61\% از حکم‌های «توانی» \lr{BIC} از داده‌های تولیدشده با مدل
توانی آمده‌اند؛ 38\% از \lr{expasym.)}

\subsubsection{\texorpdfstring{بخش \lr{E:} نتیجه‌گیری درباره‌ی داده‌های
هفته‌ی 2 (15
دقیقه)}{بخش E: نتیجه‌گیری درباره‌ی داده‌های هفته‌ی 2 (15 دقیقه)}}

یک پاراگراف (6 تا 10 جمله) بنویسید که یک داور بپذیرد و خلاصه کند داده‌های
هفته‌ی 2 درباره‌ی شکل فراموشی چه می‌گویند. از بخش‌های \lr{A} تا \lr{D} و
نتایج \lr{CV} هفته‌ی 6 استفاده کنید و بگویید با این طرح آزمایش چه چیزی را
\emph{نمی‌شود} نتیجه گرفت.

\subsubsection{تمرین‌ها (به ترتیب سختی)}

\textbf{تمرین 1 \ensuremath{\star}.} وزن‌های \lr{AIC} و \lr{BIC} مثال حل‌شده را با دست
(ماشین‌حساب) حساب کنید و با بخش \lr{B} مقایسه کنید.

\textbf{تمرین 2 \ensuremath{\star}.} بخش \lr{B} را با \lr{AICc} تکرار کنید. برای چه
\(N\)ای \lr{AICc} و \lr{AIC} برای مدل دوردی به‌طور قابل توجهی فرق می‌کنند؟

\textbf{تمرین 3 \ensuremath{\star}\ensuremath{\star} (طراحی برای بازیابی).} بازیابی \lr{expasym} بخش
\lr{C} را با فاصله‌های گسترش‌یافته به
\(\{0.25,\allowbreak 1,\allowbreak 2,\allowbreak 4,\allowbreak 7,\allowbreak 14,\allowbreak 30,\allowbreak 60,\allowbreak 90\}\) روز، و جداگانه با \(n = 200\)
آیتم، دوباره اجرا کنید. کدام تغییر بیشتر به بازیابی \(\lambda\) و \(s\)
کمک می‌کند؟

\textbf{تمرین 4 \ensuremath{\star}\ensuremath{\star} \lr{(CV} در بازیابی مدل).} در بخش \lr{D}،
\lr{AIC/BIC} را با \lr{CV} پنج‌بخشی (از هفته‌ی 6) برای 30 شبیه‌سازی از هر
مدل جایگزین کنید. آیا \lr{CV} در بازیابی مدل‌ها از \lr{AIC} بهتر است یا
بدتر؟ آیا به زمان محاسبه‌اش می‌ارزد؟

\textbf{تمرین 5 \ensuremath{\star}\ensuremath{\star}\ensuremath{\star} (پیشرفته: بازیابی در \lr{SDT).}} برای آزمایش پاداش
هفته‌ی 4 (سه شرط، 80 + 80 \lr{trial} در هر شرط)، یک بازیابی مدل برای
«\(d'\) مشترک» در برابر «\(d'\) جداگانه» با آزمون نسبت درست‌نمایی در
\(\alpha = .05\) و با \lr{BIC} اجرا کنید. هر روش چند بار به‌اشتباه ادعا
می‌کند که پاداش حساسیت را تغییر می‌دهد؟ برای اینکه یک تفاوت واقعی حساسیت
به اندازه‌ی 0.3 را در 80\% شبیه‌سازی‌ها تشخیص دهید، در هر شرط چند
\lr{trial} لازم است؟

\subsubsection{\texorpdfstring{\lr{lab check-off}}{lab check-off}}

نشان دهید: (1) جدول معیارهای اطلاعاتی با وزن‌ها (بخش \lr{B)}، (2)
نمودارهای پراکنش بازیابی با یک جمله درباره‌ی اینکه کدام پارامتر بدتر
بازیابی می‌شود و چرا، (3) ماتریس‌های درهم‌ریختگی \lr{AIC} و \lr{BIC}، و (4)
پاراگراف بخش \lr{E.}

\begin{center}\rule{0.5\linewidth}{0.5pt}\end{center}

\subsection{\texorpdfstring{ارتباط با هوش مصنوعی و
\lr{LLM}ها}{ارتباط با هوش مصنوعی و LLMها}}

\textbf{مقایسه‌ی منحنی‌های یادگیری یک مدل هوش مصنوعی.} \lr{LLM}ها می‌توانند
یک تکلیف تازه را از چند مثال در \lr{prompt} یاد بگیرند («یادگیری
درون‌بافتی» یا \lr{in-context learning).} اینکه خطا با تعداد مثال‌ها \(k\)
چطور کم می‌شود یک منحنی یادگیری است، و سؤال «نمایی یا توانی؟» همان سؤال
\lr{lab} این هفته (و منحنی‌های فراموشی هفته‌ی 2) است. معیارهای اطلاعاتی،
بازیابی مدل و \lr{identifiability} بدون تغییر به کار می‌آیند. طرح هم به
همان اندازه مهم است: با مقادیر کم \(k\) و آیتم‌های آزمون کم، ممکن است دو
منحنی اصلاً قابل تمیز نباشند.

\begin{quote}
\textbf{\lr{Code} · کد:} \lr{\nolinkurl{code/week07.py}}, \lr{cell} 6 ·
سلول 6
\end{quote}

\textbf{خروجی مورد انتظار:}

\Needspace{6\baselineskip}
\begin{latin}
\begin{Verbatim}
{'exponential': 41.6, 'power': 40.4} Akaike weights: [0.362 0.638]
\end{Verbatim}
\end{latin}

داده از یک قانون توانی تولید شده، ولی شواهد ضعیف است (وزن‌های 0.64 در
برابر 0.36).

\textbf{امتحانش کنید (بدون نیاز به \lr{API).}} - \ensuremath{\star} در همین طرح یک
بازیابی مدل (هفته‌ی 7، بخش \lr{D)} اجرا کنید: 50 مجموعه داده از هر منحنی.
\lr{confusion matrix} را بسازید. - \ensuremath{\star}\ensuremath{\star} طرح را عوض کنید: مقادیر بیشتری
برای \(k\) (تا 128) یا 200 آیتم آزمون برای هر \(k\). کدام تغییر بازیابی
مدل را بیشتر بهتر می‌کند؟ - \ensuremath{\star}\ensuremath{\star} بحث: مقاله‌ای از یک منحنی برازش‌شده نتیجه
می‌گیرد که «یادگیری درون‌بافتی، مثل تمرین در انسان، از قانون توانی پیروی
می‌کند». پیش از پذیرفتن این مقایسه با انسان، چه چیزی باید ببینید؟

\begin{center}\rule{0.5\linewidth}{0.5pt}\end{center}

\subsection{اشتباهات رایج و برداشت‌های نادرست}

\begin{enumerate}
\def\labelenumi{\arabic{enumi}.}
\tightlist
\item
  \textbf{تفسیر یک مقدار \lr{AIC} به‌تنهایی.} فقط تفاوت بین مدل‌هایی که
  روی داده‌ی یکسان برازش شده‌اند معنا دارد.
\item
  \textbf{استفاده از \lr{LRT} برای مدل‌های \lr{non-nested}،} یا خواندن یک
  \lr{LRT} غیرمعنادار به‌عنوان شاهدی \emph{به نفع} مدل ساده‌تر.
\item
  \textbf{فکر کردن اینکه \lr{BIC} همان \lr{«AIC} بیزی» با همان هدف است.}
  \lr{BIC} به دنبال شناسایی مدل تولیدکننده‌ی داده است؛ \lr{AIC} به دنبال
  پیش‌بینی. می‌توانند به دلایل موجه با هم مخالف باشند.
\item
  \textbf{\(N\) نامشخص.} همیشه گزارش کنید چه چیزی به‌عنوان یک مشاهده
  شمرده شده؛ \lr{BIC} و \lr{AICc} به آن بستگی دارند.
\item
  \textbf{نادیده گرفتن بازیابی.} پارامتری که با طرح آزمایش شما قابل
  بازیابی نیست نباید تفسیر شود، هر قدر هم برازش خوب باشد.
\item
  \textbf{اجرای بازیابی با طرح یا رویه‌ی برازشی متفاوت} از تحلیل واقعی.
  نتایج بازیابی وابسته به طرح آزمایش‌اند.
\item
  \textbf{خواندن \lr{confusion matrix} به جای \lr{inversion matrix.}}
  «داده‌های مدل \lr{A} در 90\% موارد به‌عنوان \lr{A} بازیابی می‌شوند» با
  «وقتی \lr{A} برنده می‌شود در 90\% موارد مدل واقعی است» یکی نیست.
\item
  \textbf{انتخاب دامنه‌های پارامتر برای بازیابی که از مقادیر معقول پهن‌تر
  یا باریک‌ترند.} دامنه‌ی خیلی پهن همبستگی‌ها را باد می‌کند؛ دامنه‌ی خیلی
  باریک مشکلات را پنهان می‌کند.
\end{enumerate}

\begin{center}\rule{0.5\linewidth}{0.5pt}\end{center}

\subsection{خودآزمایی}

\textbf{1.} فرمول‌های \lr{AIC} و \lr{BIC} را بنویسید. برای چه \(N\)ای
\lr{BIC} هر پارامتر را بیشتر از \lr{AIC} جریمه می‌کند؟ \textbf{پاسخ}
\(\mathrm{AIC} = -2\hat\ell + 2k\)؛
\(\mathrm{BIC} = -2\hat\ell + k\ln N\). جریمه‌ی \lr{BIC} وقتی از \lr{AIC}
بیشتر است که \(\ln N > 2\)، یعنی \(N > e^2 \approx 7.4\)، پس برای
\(N \ge 8\).

\textbf{2.} مدل \lr{A} دارای \lr{AIC} = 210.4 و مدل \lr{B} دارای
\lr{AIC} = 212.9 است. وزن‌های آکائیکه را حساب و تفسیر کنید. \textbf{پاسخ}
\(\Delta = (0, 2.5)\)، پس \(e^{-\Delta/2} = (1, 0.287)\) و
\(w = (0.777, 0.223)\). \lr{A} حدود 3.5 به 1 به‌عنوان پیش‌بینی‌کننده‌ی بهتر
ترجیح دارد، ولی \lr{B} هنوز پشتیبانی قابل توجهی دارد.

\textbf{3.} چرا آزمون نسبت درست‌نمایی می‌تواند نمایی و نمایی با مجانب را
مقایسه کند ولی نمایی و توانی را نه؟ \textbf{پاسخ} نمایی حالت خاصی از مدل
مجانبی است (\(s = 0\))، پس مدل‌ها \lr{nested} هستند و \(2\Delta\ell\) تحت
فرض صفر توزیع \(\chi^2\) معلومی دارد. نمایی و توانی \lr{nested} نیستند:
هیچ‌کدام حالت خاص دیگری نیست.

\textbf{4.} مراحل یک مطالعه‌ی بازیابی پارامتر را شرح دهید. \textbf{پاسخ}
مقادیر واقعی معقول برای پارامترها انتخاب کنید؛ با طرح واقعی داده
شبیه‌سازی کنید؛ با رویه‌ی برازش واقعی برازش کنید؛ مقادیر بازیابی‌شده را با
واقعی مقایسه کنید (نمودار پراکنش، همبستگی، سوگیری، \lr{RMSE)}؛ همبستگی
بین پارامترهای بازیابی‌شده را برای یافتن \lr{trade-off} بررسی کنید.

\textbf{5.} در یک مطالعه‌ی بازیابی، \(\lambda\) و \(s\) بازیابی‌شده با
\(r = -0.8\) همبسته‌اند، با اینکه مقادیر واقعی مستقل کشیده شده بودند. این
یعنی چه؟ \textbf{پاسخ} یک \lr{trade-off:} با این طرح آزمایش، ترکیب‌های
مختلف نرخ زوال و مجانب پیش‌بینی‌های تقریباً یکسانی می‌دهند، پس خطا در یکی با
دیگری جبران می‌شود. مقادیر تکی \(\lambda\) و \(s\) نباید جداگانه تفسیر
شوند. طرح آزمایش بهتر یا \lr{reparameterization} لازم است.

\textbf{6.} تفاوت \lr{confusion matrix} و \lr{inversion matrix} را در
بازیابی مدل توضیح دهید. \textbf{پاسخ} درهم‌ریختگی:
\(P(\text{\lr{selected}} \mid \text{\lr{generated}})\) که سطرهایش جمع 1 دارند و
رفتار روش را برای هر مدل واقعی توصیف می‌کند. معکوس:
\(P(\text{\lr{generated}} \mid \text{\lr{selected}})\) که ستون‌هایش جمع 1 دارند (با
فرض پیشین برابر برای مدل‌های تولیدکننده) و به ما می‌گوید هنگام تحلیل داده‌ی
واقعی چقدر به یک برنده‌ی مشخص اعتماد کنیم.

\textbf{7.} یک بازیابی مدل نشان می‌دهد وقتی مدل دوپارامتری داده را تولید
کرده، \lr{AIC} در 40\% موارد مدل سه‌پارامتری را انتخاب می‌کند. چه می‌کنید؟
\textbf{پاسخ} با این طرح آزمایش، به برد \lr{AIC} برای مدل پیچیده با
احتیاط نگاه کنید: بازیابی \lr{BIC} یا \lr{CV} را بررسی کنید که ممکن است
محافظه‌کارتر باشند، طرح آزمایش را بهتر کنید \lr{(trial} بیشتر، شرایط
آگاهی‌بخش‌تر)، و نتایج بازیابی را کنار مقایسه‌ی داده‌ی واقعی گزارش کنید.

\begin{center}\rule{0.5\linewidth}{0.5pt}\end{center}

\subsection{تکلیف}

\subsubsection{\texorpdfstring{ادامه‌ی \lr{A3} (موعد ابتدای درس هفته‌ی
9)}{ادامه‌ی A3 (موعد ابتدای درس هفته‌ی 9)}}

کارهای 4 (معیارهای اطلاعاتی) و 5 (بازیابی پارامتر و مدل) \lr{A3} از
مطالب این هفته استفاده می‌کنند. شرح کامل در هفته‌ی 6 است. کد
\lr{\nolinkurl{parameter_recovery}} امروز و حلقه‌ی بازیابی مدل را دوباره
به کار ببرید: طوری طراحی شده‌اند که مستقیماً قابل تطبیق باشند.

\subsubsection{پروژه‌ی پایانی: تشکیل گروه‌ها و انتشار راهنمای پیشنهاد}

\begin{itemize}
\tightlist
\item
  \textbf{گروه‌ها:} 1 یا 2 نفر. تا پایان این هفته گروهتان را نزد دستیار
  آموزشی ثبت کنید.
\item
  \textbf{کار:} یک مدل شناختی انتخاب کنید و کل مسیر مدل‌سازی را روی یک
  مجموعه‌داده (شبیه‌سازی‌شده، باز، یا داده‌ی خودتان) اجرا کنید: برازش،
  ارزیابی (دست‌کم یک مقایسه با یک مدل جایگزین)، بازیابی (دست‌کم یک تحلیل
  بازیابی پارامتر)، تفسیر. \textbf{گروه‌های دونفره} یک تحلیل
  بیزی/\lr{MCMC} یا سلسله‌مراتبی و یک بیانیه‌ی سهم فردی هم اضافه می‌کنند.
\item
  \textbf{پیشنهاد پروژه (موعد ابتدای درس هفته‌ی 9، حداکثر 2 صفحه):} سؤال
  پژوهش، مدل(ها) با معادلات، منبع داده و طرح آزمایش، تحلیل‌های مقایسه و
  بازیابی برنامه‌ریزی‌شده، و زمان‌بندی. شرح کامل پروژه و روبریک جداگانه
  توزیع می‌شود (\lr{\nolinkurl{project_brief.md}}).
\item
  \textbf{مطمئن نیستید چه چیزی را مدل کنید؟} موضوعات هفته‌های 14\lr{--}15
  (دسته‌بندی \lr{exemplar}، \lr{drift-diffusion}، یادگیری تقویتی) همه
  پروژه‌های خوبی هستند، همین‌طور گسترش مدل‌های فراموشی و \lr{SDT} این درس.
\end{itemize}

\begin{center}\rule{0.5\linewidth}{0.5pt}\end{center}

\subsection{مطالعه‌ی بیشتر}

\begin{itemize}
\tightlist
\item
  \textbf{\lr{Farrell} \lr{\&} \lr{Lewandowsky} (2018)، فصل 10.} مقایسه‌ی
  مدل‌ها: آزمون نسبت درست‌نمایی، \lr{AIC}، \lr{BIC} و وزن‌هایشان؛ فرض‌ها و
  محدودیت‌هایشان. \emph{(منبع اصلی.)}
\item
  \textbf{\lr{Lee} \lr{\&} \lr{Wagenmakers} (2014)، فصل 7} (مقدمه). اینکه
  رویکرد بیزی به مقایسه‌ی مدل‌ها چه ارتباطی با روش‌های این هفته دارد؛
  آمادگی برای هفته‌ی 10.
\item
  \emph{اختیاری، خارج از فهرست درس:} \lr{Wilson, R. C.,} \lr{\&} \lr{Collins,
  A. G. E.} (2019). \lr{Ten simple rules for the computational modeling
  of behavioral data.} \emph{\lr{eLife,} 8}, \lr{e49547.} راهنمایی عملی
  که بازیابی پارامتر و مدل را رویه‌ی استاندارد می‌داند.
\end{itemize}

\emph{پیش‌نمایش هفته‌ی 8:} میان‌ترم. درس یک مرور یکپارچه است که در آن کل
مسیر (شبیه‌سازی، برازش، ارزیابی، بازیابی، تفسیر) را روی یک مجموعه‌داده‌ی
جدید اجرا می‌کنیم؛ زمان \lr{lab} به آزمون میان‌ترم اختصاص دارد.
